\section{Proof of the all-time theorem}
\label{app:proof-alltime}

We prove \Cref{thm:alltime} for layer-dependent activations.  Write
$g_\ell=(\phi_\ell)'$, and fix bounds
\[
 |\phi_\ell(0)|\le a,\qquad |g_\ell|\le s,\qquad
 |g_\ell(u)-g_\ell(v)|\le j|u-v|.
\]
All constants below may depend on these bounds, the fixed data and depth,
and the initial Gram margin.  They do not depend on $n,q$, or time.
Constants denoted by $C$ may increase from line to line.  For sample vectors
write $\|v\|_m^2=m^{-1}\sum_a v_a^2$, so $Y=\|y\|_m$.
Vector and matrix norms at finite width are ordinary Euclidean and
Frobenius norms; every normalization by $n$ is displayed.
The distance $d_n$ is the sum of scaled block norms in
\eqref{eq:normalized-distance}.  We first prove uniform fitting at every
order, then construct the dense Gaussian reference and transfer its carrier
tails, and finally use these tails to sharpen the closure error.

\subsection{Initialization and polynomial projection}

Let $v_a=x_a/\sqrt d$, $X=\max_a\|v_a\|_2$, and define the deterministic
uncentered feature covariances
\begin{align}
 Q^{(1)}_{ab}&=\E[\phi_1(G^\top v_a)\phi_1(G^\top v_b)],
       &&G\sim N(0,I_d),\nonumber\\
 Q^{(\ell)}_{ab}&=\E[\phi_\ell(Z_a)\phi_\ell(Z_b)],
       &&Z\sim N(0,Q^{(\ell-1)}),\quad\ell\ge2.
 \label{eq:at-initial-covariance}
\end{align}
Choose $\lambda>0$ with $Q^{(L)}/m\succeq2\lambda I_m$.
There are fixed $K_0,A_0,C_0$ such that the events
\begin{equation}
 \mathcal G_n=\left\{\begin{aligned}
 &
 \max_{\ell\ge2}\|W_0^{(\ell)}\|_{\rm op}\le K_0,\quad
 \max_a\frac{\|z_a^{(1)}(0)\|_2}{\sqrt n}\le A_0,\\
 &
 \frac{\|W_0^{(1)}\|_F}{\sqrt n}\le C_0,\quad
 \Gamma_w(0):=\frac{H_L(0)^\top H_L(0)}{mn}\succeq\lambda I_m
 \end{aligned}\right\}
 \label{eq:at-good-event}
\end{equation}
have probability tending to one.  Here $H_L$ has feature columns indexed
by the training samples.  Indeed, a $1/4$-net of the unit sphere has at
most $9^n$ points, and approximation of both vectors in a bilinear form
gives
\[
 \Pr\{\|W_0^{(\ell)}\|_{\rm op}>K_0\}
 \le2\,9^{2n}e^{-nK_0^2/8}\longrightarrow0
\]
for a sufficiently large fixed $K_0$.  The first-layer bounds follow from
the law of large numbers at fixed $d,m$.

For completeness, $H_\ell(0)^\top H_\ell(0)/n\to Q^{(\ell)}$ in
probability.  At the first layer this is a law of large numbers.
Conditionally on all previous layers, the rows of the next preactivation
matrix are independent $N(0,H_{\ell-1}^\top H_{\ell-1}/n)$ vectors.
Linear growth, $|\phi(u)|\le a+s|u|$, bounds each activation's Gaussian
fourth moment by $8a^4+24s^4Q_{aa}^2$.  Conditional Chebyshev therefore
bounds each covariance-entry error by $C(M)/(n\epsilon^2)$ whenever the
previous covariance has operator norm at most $M$.  Finally the map
$Q\mapsto\E[\phi(Q^{1/2}G)_a\phi(Q^{1/2}G)_b]$ is continuous,
including at singular $Q$: positive square roots converge, and a constant
multiple of $1+\|G\|_2^2$ is an integrable majorant on bounded covariance
sets.  Induction proves the assertion and hence the Gram part of
\eqref{eq:at-good-event}.

The input criterion stated with the theorem also follows directly.
If the nonzero $v_a$ are pairwise nonproportional and
$\sum_a c_a\phi_1(v_a^\top u)=0$ for all $u$, fix $a$ with
$c_a\ne0$.  For each $b\ne a$, choose $u_b\perp v_b$ with
$u_b\cdot v_a\ne0$.  Applying all $m-1$ finite differences in the
directions $t_bu_b$ eliminates every term but the $a$th, giving
\[
 \Delta_{s_1}\cdots\Delta_{s_{m-1}}\phi_1(t)=0
 \quad\text{for all }t,s_1,\ldots,s_{m-1}.
\]
Convolve with a smooth compactly supported approximate identity and
differentiate once in each $s_i$ at zero.  The mollified function is a
polynomial of degree at most $m-2$.  Its coefficients converge, by
interpolation at $m-1$ distinct points, so the original continuous
function is such a polynomial.  Thus a nonpolynomial first activation
makes the ridge functions independent.  A null quadratic form of
$Q^{(1)}$ would give a continuous ridge combination vanishing almost
everywhere under a full-support Gaussian, hence everywhere, and is
impossible.  If $Q\succ0$ and $\psi$ is continuous and nonconstant,
the identity $\sum_a c_a\psi(Z_a)=0$ for $Z\sim N(0,Q)$ likewise holds
everywhere; varying one coordinate shows $c_a=0$.  This propagates
positivity to all layers.  For $m=1$, a nonzero input and any nonconstant
first activation suffice.  Alternatively linearly independent inputs
give a positive first preactivation covariance, so any nonconstant
activation in every layer suffices.  A globally Lipschitz nonaffine
function is nonpolynomial.  These are sufficient criteria; the proof
below only needs the actual gap in \eqref{eq:at-initial-covariance}.

\begin{lemma}[Projection estimates used in the bootstrap]
\label{lem:at-projection}
Let $\Pi_q^A$ be orthogonal projection onto polynomials of degree below
$q$ in $L^2(0,A;\mathcal H)$, for a real Hilbert space $\mathcal H$.
For $u\in H^1(0,A;\mathcal H)$ and bounded measurable $b$,
\begin{align}
 \|(I-\Pi_q^A)u\|_{L^2}^2
 &\le\frac1{q(q+1)}\int_0^A\xi(A-\xi)\|u'(\xi)\|^2\dd\xi,
 \label{eq:at-weighted-tail}\\
 \|u(A)-(\Pi_q^Au)(A)\|
 &\le C\sqrt{A/q}\,\|u'\|_{L^2},&
 \|(\Pi_q^Ab)(A)\|&\le C\sqrt q\,\|b\|_{L^\infty}.
 \label{eq:at-endpoint}
\end{align}
If a history $u$ is fixed as its right endpoint grows, then
\begin{equation}
 \frac{\dd}{\dd A}\|(I-\Pi_q^A)u\|_{L^2(0,A)}^2
   =\|u(A)-(\Pi_q^Au)(A)\|^2
 \quad\text{for almost every }A.
 \label{eq:at-growing-energy}
\end{equation}
\end{lemma}
\begin{proof}
The normalized shifted Legendre modes $e_k$ satisfy
$-[\xi(A-\xi)e_k']'=k(k+1)e_k$.
Integration by parts and Bessel's inequality applied to the orthonormal
family $\sqrt{\xi(A-\xi)}e_k'/\sqrt{k(k+1)}$ give
\eqref{eq:at-weighted-tail}; the omitted modes have $k(k+1)\ge q(q+1)$.
Smooth approximation and an orthonormal expansion in $\mathcal H$ give
the stated generality.

On $[0,1]$ the endpoint kernel is
$K_q(u)=\sum_{k<q}(2k+1)p_k(u)=(p_q'(u)+p_{q-1}'(u))/2$.
Integration by parts yields
\[
 u(A)-(\Pi_q^Au)(A)
 =\frac12\int_0^A[p_q(\xi/A)+p_{q-1}(\xi/A)]u'(\xi)\dd\xi.
\]
The squared $L^2$ norm of this multiplier is
$A[(2q+1)^{-1}+(2q-1)^{-1}]/4$, proving the first endpoint bound.
For the other one it suffices to check $\int_0^1|K_q|\le C\sqrt q$.
Here is a short verification of the needed Legendre variation estimate.
For the ordinary Legendre polynomial $L_k$, set
$v(\theta)=\sqrt{\sin\theta}L_k(\cos\theta)$ and
$Q(\theta)=(k+1/2)^2+(4\sin^2\theta)^{-1}$.
The Legendre equation gives $v''+Qv=0$ and
$(v^2+(v')^2/Q)'=-Q'(v')^2/Q^2\ge0$ on $(0,\pi/2)$.
The values $L_{2r}(0)=(-1)^r\binom{2r}{r}/4^r$,
$L_{2r+1}(0)=0$, and $L_k'(0)=kL_{k-1}(0)$ show that this energy
at $\pi/2$ is at most $C/k$, using
$\binom{2r}{r}/4^r\le(r+1)^{-1/2}$.
Consequently, on $1/k\le\theta\le\pi/2$,
\[
 |\partial_\theta L_k(\cos\theta)|
 \le Ck^{-1/2}\{k(\sin\theta)^{-1/2}+(\sin\theta)^{-3/2}\}.
\]
Its integral is at most $C\sqrt k$.  On $[0,1/k]$, use
$|L_k'|\le k(k+1)/2$; this follows by summing
$L_{r+1}'-L_{r-1}'=(2r+1)L_r$ and $|L_r|\le1$ on $[-1,1]$.
The last bound follows, for example, by integrating
$(\cos\theta+i\sin\theta\cos\alpha)^r$ over $0\le\alpha\le\pi$.
Parity handles the other half.  Thus ${\rm Var}(L_k)\le C\sqrt k$,
which proves the kernel bound; $q=1$ is immediate.

Finally differentiate the minimum least-squares error on $[0,A]$.
The endpoint contributes the right side of \eqref{eq:at-growing-energy}.
The other term pairs the projection residual with the derivative of a
polynomial of degree below $q$ and vanishes by orthogonality.  Moment
coordinates and the invertible polynomial Gram matrix justify absolute
continuity of its coefficients on compact endpoint intervals.  This
argument only needs square-integrable histories for the last identity.
\end{proof}

\subsection{Fitting and continuation for every memory order}
\label{sec:at-fitting}

Fix an initialization in $\mathcal G_n$ and put $D_0=K_0+1$.
On the tube $\|W^{(\ell)}\|_{\rm op}<D_0$ for $\ell\ge2$ and
$\max_a\|z_a^{(1)}\|_2/\sqrt n<A_0+1$, define
\[
 M_1=a+s(A_0+1),\qquad M_\ell=a+sD_0M_{\ell-1}.
\]
Linear growth gives $\|h_a^{(\ell)}\|_2/\sqrt n\le M_\ell$.
Stop the local closure at first tube exit or activity
$S=4Y/\lambda$, reducing the eventual $Y_*$ so $S\le1$.
Since $w(0)=0$, the readout equation and backward recursion give
\begin{equation}
 \frac{\|w\|_2}{\sqrt n}\le2M_LS,
 \qquad \max_a\frac{\|\delta_a^{(\ell)}\|_2}{\sqrt n}
 \le\beta_\ell,\qquad
 \beta_L=2sM_LS,\quad\beta_\ell=sD_0\beta_{\ell+1}.
 \label{eq:at-rms}
\end{equation}
Thus every $\beta_\ell\le CY$.  In this subsection responses and
residuals belong to the closure, with hats suppressed.
Write $c_a=r_a/\rho$, $b_a^{(\ell)}=c_a\delta_a^{(\ell)}$,
and $A=\tau(t)\le1+S\le2$.  The forward prefix is constant and the
backward prefix is zero.  Projection contraction in the reconstruction gives
\begin{align}
 \|\widehat W^{(\ell)}-W_0^{(\ell)}\|_F
 &\le2M_{\ell-1}\beta_\ell\sqrt{(1+S)S}=O(Y^{3/2}),
 \label{eq:at-hidden-tube}\\
 \frac{\|\widehat W^{(1)}-W_0^{(1)}\|_F}{\sqrt n}
 &\le2X\beta_1S=O(Y^2).
 \label{eq:at-first-tube}
\end{align}
For the first inequality, the backward history's averaged squared norm
is at most $S\beta_\ell^2$, the forward one at most
$(1+S)M_{\ell-1}^2$, and
$\|uv^\top/n\|_F=(\|u\|_2/\sqrt n)(\|v\|_2/\sqrt n)$.
These bounds keep both tube boundaries strictly away for small $Y$.

Differentiating the reconstruction gives exactly
\begin{equation}
 \dot{\widehat\theta}=F(\widehat\theta)+E,\qquad E_1=E_w=0,
 \qquad E_\ell=\frac{2\rho}{mn}\sum_a
       (b_a^{(\ell)}-b_a^{(\ell)*})
       (h_a^{(\ell-1)}-h_a^{(\ell-1)*})^\top,
 \label{eq:at-defect}
\end{equation}
where stars denote projected endpoint values.  By
\eqref{eq:at-growing-energy}, zero initial projection errors, and
Cauchy--Schwarz, its accumulated absolute norm satisfies
\begin{equation}
 \int_0^t\|E_\ell(u)\|_F\dd u
 \le\frac2{mn}\sum_a
 \|(I-\Pi_q^A)b_a^{(\ell)}\|_{L^2(0,A)}
 \|(I-\Pi_q^A)h_a^{(\ell-1)}\|_{L^2(0,A)}.
 \label{eq:at-defect-product}
\end{equation}
This bounds total absolute velocity error, not just the error in a signed
history integral.

Set
\[
 Z_\ell(t)=\frac1{mn}\sum_a\int_0^{\tau(t)}
       \|\partial_\xi h_a^{(\ell)}\|_2^2\dd\xi.
\]
First $Z_1\le4Ss^2X^4\beta_1^2$.
Endpoint evaluation on degree-below-$q$ polynomials has norm
$q/\sqrt A$ because $\sum_{k<q}(2k+1)=q^2$.
The averaged backward endpoint error is consequently at most
$(q+1)\beta_\ell$ in normalized Euclidean norm.  Equations
\eqref{eq:at-weighted-tail}--\eqref{eq:at-defect} imply
\[
 \int_0^t\rho\|E_\ell/\rho\|_F^2\dd u
 \le2(1+S)^2\beta_\ell^2 Z_{\ell-1}(t).
\]
The clock chain rule is
\[
 \partial_\xi h_a^{(\ell)}=g_\ell(z_a^{(\ell)})\od
 \{(F_\ell/\rho+E_\ell/\rho)h_a^{(\ell-1)}
                   +\widehat W^{(\ell)}\partial_\xi h_a^{(\ell-1)}\}.
\]
As $\|F_\ell/\rho\|_F\le2\beta_\ell M_{\ell-1}$, the preceding
estimate proves inductively
\begin{align}
 Z_1&\le Z_1^*:=4Ss^2X^4\beta_1^2,\nonumber\\
 Z_\ell&\le Z_\ell^*:=3s^2\left[
 4S\beta_\ell^2M_{\ell-1}^4+
 \{D_0^2+2(1+S)^2\beta_\ell^2M_{\ell-1}^2\}Z_{\ell-1}^*\right]
 \le CY^3.
 \label{eq:at-forward-energy}
\end{align}
The factors $M_{\ell-1}^4$ and $M_{\ell-1}^2$ retain the contribution
of unbounded activation values.  No feature coordinate maximum is used.
The sharper endpoint bounds \eqref{eq:at-endpoint} now give forward
error at most $CY^{3/2}/\sqrt q$ and backward error at most
$CY\sqrt q$.  Their product in \eqref{eq:at-defect} yields
\begin{equation}
 e_E(t):=\sum_{\ell=2}^L\|E_\ell(t)\|_F
 \le CY^{5/2}\rho(t),\qquad
 \int_0^t e_E\dd u\le CY^3/q.
 \label{eq:at-relative-defect}
\end{equation}
The second estimate also follows from \eqref{eq:at-defect-product},
the $C Y^{3/2}/q$ forward tail and the $C Y^{3/2}$ backward history norm.

The tangent Gram in the residual equation is
\begin{equation}
 \Gamma_{ab}=\frac1m\left[
 \frac{(h_a^{(L)})^\top h_b^{(L)}}n+
 \frac{(\delta_a^{(1)})^\top\delta_b^{(1)}}n\,v_a^\top v_b+
 \sum_{\ell=2}^L
 \frac{(\delta_a^{(\ell)})^\top\delta_b^{(\ell)}}n
 \frac{(h_a^{(\ell-1)})^\top h_b^{(\ell-1)}}n\right].
 \label{eq:at-tangent-gram}
\end{equation}
Every summand is a parameter-gradient Gram, so $\Gamma\succeq\Gamma_w$.
Equation \eqref{eq:at-forward-energy} gives
\[
 \frac{\|H_L(t)-H_L(0)\|_F}{\sqrt{mn}}\le\sqrt{SZ_L^*},\qquad
 \|\Gamma_w(t)-\Gamma_w(0)\|_{\rm op}
       \le2M_L\sqrt{SZ_L^*}\le CY^2.
\]
After reducing $Y_*$, $\Gamma\succeq\lambda I/2$ and
$\|\Gamma\|_{\rm op}\le G_*$ on the stopped interval.
For the ordinary prediction Jacobian $J$,
$(JE)_a=\sum_{\ell\ge2}(\delta_a^{(\ell)})^\top E_\ell
h_a^{(\ell-1)}/n$, whence
\[
 \dot r=-2\Gamma r+JE,\qquad
 \|JE\|_m\le CY e_E\le CY^{7/2}\rho.
\]
Choose $Y_*$ so $CY_*^{7/2}\le\lambda/2$ and put $\kappa=\lambda/2$.
Taking the inner product with $r/\rho$ proves
\begin{equation}
 Y e^{-\Lambda t}\le\rho(t)\le Ye^{-\kappa t},\qquad
 \int_0^t\rho\dd u\le Y/\kappa=S/2,
 \label{eq:at-fitting}
\end{equation}
for a fixed finite $\Lambda$.  Thus the activity boundary is also excluded.

These arguments legitimately use the clock before its lower residual bound
has been proved.  The raw moment vector field is locally Lipschitz for
$\tau>0$, and all its velocities vanish at zero residual.  Local
uniqueness backwards from such an equilibrium excludes a first finite zero
on a nonstationary regular solution.  At fixed $n,q$, the physical bounds
above and the integral moment formulas bound the entire raw state in a
compact subset of $\tau>0$.  A finite maximal endpoint would therefore
extend.  This proves global existence, uniqueness and \eqref{eq:at-fitting}
for every $q\ge1$.  The dense flow has the same proof with $E=0$ and
hidden displacement at most $2S\beta_\ell M_{\ell-1}$.
All reductions of $Y_*$ impose finitely many inequalities whose left
sides vanish with $Y$; none involves width or order.

The update equations, \eqref{eq:at-relative-defect}, and forward chain rule
also give, for either path with its own residual,
\begin{align}
 \frac{\|\dot W^{(1)}\|_F}{\sqrt n}
  +\sum_{\ell=2}^L\|\dot W^{(\ell)}\|_F&\le CY\rho,
 &\frac{\|\dot w\|_2}{\sqrt n}&\le C\rho,\nonumber\\
 \max_{\ell,a}\frac{\|\dot z_a^{(\ell)}\|_2+
                      \|\dot h_a^{(\ell)}\|_2}{\sqrt n}&\le CY\rho,
 &\int_t^\infty\rho(u)\dd u&\le\rho(t)/\kappa.
 \label{eq:at-speeds}
\end{align}
Consequently physical parameters converge to finite interpolating limits.
The residual equation further gives
$\|\dot r\|_m\le C\rho$ and $\|\dot c\|_m\le C$.
When $Y=0$, both paths are stationary and all conclusions are immediate.

\subsection{The Gaussian reference and its finite-width transfer}
\label{sec:at-gaussian}

We give the Gaussian construction needed below, including the fixed
computation input.  Its purpose is to control full dense backward carriers,
including the trained readout.  All Gaussian bounds in this subsection
initially concern population reference computations; no Gaussian law for
the coordinates of a trained finite network is assumed.

\begin{lemma}[Fixed Gaussian computations]
\label{lem:at-fixed-program}
Consider a fixed finite list of typed vector instructions using independent
matrices with entries $N(0,1/n)$, both orientations of each same matrix,
independent identically distributed root tuples in each layer, deterministic
linear combinations, and $C^1$ coordinate maps with bounded first derivatives.
Root tuples in different layers are independent and have finite second
moments.  Every same-layer tuple has a deterministic limiting empirical law,
in probability in Wasserstein distance of order two.  Finite collections
of such programs sharing the arrays converge jointly.

The limiting forward and reverse actions of one initialized matrix are
\begin{equation}
 \mathscr W h=\xi_h+\sum_{s:\,\mathscr W^*u_s\ {\rm earlier}}
       u_s\,\E[\partial_{\zeta_s}h],\qquad
 \mathscr W^*u=\zeta_u+\sum_{r:\,\mathscr Wv_r\ {\rm earlier}}
       v_r\,\E[\partial_{\xi_r}u].
 \label{eq:at-source-rule}
\end{equation}
The centered Gaussian source covariances are
$\E[\xi_h\xi_v]=\E[hv]$ and $\E[\zeta_u\zeta_v]=\E[uv]$.
Distinct matrix/orientation source groups are independent of one another
and of the roots.  Each derivative is of the unrolled scalar expression
in named source coordinates, holding deterministic coefficients,
covariances and previously computed expectations fixed.
The statement allows singular source covariances.

The value conclusion also holds for continuous coordinate maps of at most
linear growth.  If roots are sub-Gaussian, activations are $C^2$ with
bounded first two derivatives, and the extra instructions are products
$b(z)p$ with bounded $C^1$ functions $b,b'$, these products also obey the
response rule \eqref{eq:at-source-rule}.
\end{lemma}
\begin{proof}
We record the conditioning argument to specify precisely the result being
used.  For one matrix $W$, collect previous constraints as $WV=Y$ and
$W^\top U=Q$.  Conditional on the computation transcript, query inputs
are fixed.  Revealing a new answer conditions only the queried matrix's
current Gaussian factor on one more linear constraint; coordinate
operations reveal no additional randomness.  Induction therefore leaves
the residual Gaussian factors for distinct matrices conditionally
independent.  When the two query Grams are invertible, Gaussian
orthogonal projection gives
\[
 W\mid\mathcal H\ \overset d=
 Y(V^\top V)^{-1}V^\top+
 U(U^\top U)^{-1}Q^\top P_{V^\perp}
 +P_{U^\perp}\widetilde W P_{V^\perp}.
\]
Indeed the displayed mean solves both constraints, using
$U^\top Y=Q^\top V$, and is Frobenius-orthogonal to their common
homogeneous solution space.  The last term is the Gaussian on that space.
For a new input $h$, write
$\alpha_n=(V^\top V)^{-1}V^\top h$ and $h_\perp=h-V\alpha_n$.
Then
\[
 Wh=Y\alpha_n+U\beta_n+
       \frac{\|h_\perp\|_2}{\sqrt n}P_{U^\perp}g,\qquad
 \beta_n=(U^\top U/n)^{-1}(Q^\top h_\perp/n),
\]
in conditional law, with independent $g\sim N(0,I_n)$.
If the limiting query Grams are positive definite, all coefficients
converge by induction.  The removed noise has conditional expected
normalized squared norm $\operatorname{rank}(U)/n$, and hence vanishes.
For bounded Lipschitz empirical tests, the remaining independent Gaussian
coordinates give conditional variance at most $4\|\psi\|_\infty^2/n$.
The conditional mean is an empirical continuous function of the old
tuple.  Second moments converge as well: the new Gaussian cross term
has variance $4\sigma^2\|m\|_2^2/n^2$ when the conditional mean is $m$,
and the Gaussian squared-norm average has variance $2/n$.
Thus weak convergence and second-moment convergence, equivalently
Wasserstein convergence of order two, propagate through a matrix call.
Lipschitz coordinate instructions propagate them directly.

To identify the mean, let $\alpha$ be the limiting regression coefficient
and $h_\perp=h-\sum_r\alpha_rv_r$.  Earlier reverse answers equal
$\zeta_s$ plus linear combinations of the $v_r$; their pairings with
$h_\perp$ are therefore $\E[\zeta_sh_\perp]$.
Gaussian integration by parts gives
\[
 \E[\zeta h_\perp]=(\E[u_su_t])_{st}\E[\nabla_\zeta h_\perp].
\]
After inserting the earlier forward decompositions in $Y\alpha$,
the terms containing $\alpha_r\E[\nabla_\zeta v_r]$ cancel.
The remaining response is $\sum_su_s\E[\partial_{\zeta_s}h]$.
The Gaussian part has variance $\E[h^2]$ and covariance $\E[hv_r]$
with each previous forward source.  Its fresh innovation is independent
of other source groups.  This proves \eqref{eq:at-source-rule};
interchanging the two orientations proves the reverse formula.

To remove the rank assumption, add a fresh independent input
$\epsilon\chi$ at each matrix call.  At fixed $\epsilon>0$ the squared
distance of the new query from the previous query span has limiting
value at least $\epsilon^2$: the new Gaussian squared residual tends
to one, and its cross term has conditional variance $O(1/n)$ on
bounded norm events.  All limiting query Grams are then positive definite.
On the common initialized-operator event and the event
$\|\chi\|_2/\sqrt n\le2$ for the finitely many added roots, finite
instruction induction gives
$\max_v\|v_n^\epsilon-v_n\|_2/\sqrt n\le C\epsilon$.
The scalar response construction is also continuous at $\epsilon=0$.
Its covariance matrices converge, so couple their source vectors by
$K_\epsilon^{1/2}G$.  Positive square roots are continuous even at
singular matrices: bounded subsequences of the roots have limits whose
square is $K$, and uniqueness of the positive root identifies every
limit.  At a fixed finite instruction, scalar values have a uniform
linear envelope and their first source derivatives are bounded when
earlier coefficients lie in a compact set.  Dominated convergence
therefore passes values, second moments and expected source derivatives
to zero noise, closing this finite induction.  Taking width first and
then $\epsilon\downarrow0$ proves the singular case.  This argument
uses no continuity of pseudoinverses.

For a continuous map $F$ with a linear envelope, $X_j\to X$ in $L^2$
implies $F(X_j)\to F(X)$ in $L^2$: compact uniform continuity gives
convergence in probability, and the common linear envelope supplies
uniform integrability.  Approximate $F$ by smooth compactly supported
maps with a common linear envelope.  Choose each current approximation
first, then the approximation tolerance for its finite input prefix
smaller than its error budget divided by that smooth map's Lipschitz
constant.  Matrix actions multiply same-array RMS errors by at most
$K_0$.  Finite induction therefore removes these approximations in both
the scalar and empirical laws; no uniform cutoff Lipschitz constant is
needed.

For the final response assertion, replace $b(z)p$ by
$b(z)\chi_R(p)$, where $\chi_R$ is a smooth clip with
$|\chi_R(p)|\le|p|$, $|\chi_R'|\le1$, converging locally in $C^1$ to
the identity.  The clipped coordinate maps have bounded derivatives.
In a fixed scalar program the values have a linear envelope in the finite
root/source list; their first source derivatives have a polynomial
envelope, since each differentiated product adds at most one carrier
factor.  Sub-Gaussian roots and finite Gaussian source lists have all
moments.  Covariance-square-root coupling and uniform integrability thus
pass the expected derivatives through the clips, one instruction at a
time.  The value limit agrees with the preceding continuous-map
extension.  The number of instructions is fixed throughout.
\end{proof}

The same-array operator inequality used in this proof constructs actual
population operators, rather than independent reverse maps.  Include a
countable family of all required finite programs, their finite unions,
rational linear combinations, smooth coordinate approximations, and
forward and reverse queries.  Their scalar laws are consistent, since
deleting unused instructions changes no finite-array node.  Realize their
countable joint laws, separately at each layer, and complete the generated
span in $L^2$.  The high-probability inequality
$\|W_0v\|_2/\sqrt n\le K_0\|v\|_2/\sqrt n$ passes to each finite generated
span by second-moment convergence, then to its completion.  Including
a countable dense family of bounded smooth functions of finite tuples
makes this completion the $L^2$ space of the generated sigma field.
The finite identity
$n^{-1}u^\top W_0v=n^{-1}(W_0^\top u)^\top v$ passes to these spans
and then by continuity, so the reverse action is exactly $W_0^*$.
Increments of a hidden operator will be Hilbert--Schmidt:
$(u\otimes v)h=u\E[vh]$ has
$\|u\otimes v\|_{\rm HS}=\|u\|_2\|v\|_2$, the counterpart of the
finite Frobenius identity.

\paragraph{Activity-stopped population programs.}
Initially assume the activations are $C^2$ with
$\|\phi_\ell'\|_\infty\le s$ and $\|\phi_\ell''\|_\infty\le j$.
Use upper-case $Z,H,P$ for population scalar fields, with
$P_a^{(L)}=w$, $P_a^{(\ell)}=(W^{(\ell+1)})^*\delta_a^{(\ell+1)}$,
and $\delta_a^{(\ell)}=g_\ell(Z_a^{(\ell)})P_a^{(\ell)}$.
For an Euler program with physical steps $h_k$, let
\[
 \rho_k=\|r_k\|_m,\qquad \alpha_k=h_k\rho_k,\qquad
 u_{a,k}=r_{a,k}/\rho_k .
\]
Freeze zero-residual nodes.  Otherwise $|u_{a,k}|\le\sqrt m$.
The update formulas in activity are
\begin{align}
 \Delta w_k&=-\frac{2\alpha_k}{m}\sum_a u_{a,k}H_{a,k}^{(L)},\nonumber\\
 \Delta W_k^{(1)}&=-\frac{2\alpha_k}{m}\sum_a
                  u_{a,k}\delta_{a,k}^{(1)}v_a^\top,\nonumber\\
 \Delta W_k^{(\ell)}&=-\frac{2\alpha_k}{m}\sum_a
                  u_{a,k}\delta_{a,k}^{(\ell)}\otimes H_{a,k}^{(\ell-1)}.
 \label{eq:at-pop-euler}
\end{align}
At any fixed population program all coefficients here are deterministic,
although calculated causally from earlier population expectations.
Thus Lemma~\ref{lem:at-fixed-program} applies with these coefficients
frozen in every source derivative.

Stop at activity $\sum_{k<j}\alpha_k\le S\le1$ and at the physical tube
used above.  The same linear-growth induction gives $\|H_a^{(\ell)}\|_2
\le M_\ell$, $\|w\|_2\le2M_LS$, and
$\|\delta_a^{(\ell)}\|_2\le C_\ell S$.
Summing the rank-one updates gives first-layer and hidden increments
at most $CS^2$.  Subtracting successive forward recursions then gives
\[
 \max_{\ell,a}\|H_a^{(\ell)}-H_a^{(\ell)}(0)\|_2\le CS^2,\qquad
 \|\Gamma_w-\Gamma_w(0)\|_{\rm op}\le CS^2.
\]
These estimates also hold at an affine Euler interpolation state, regarded
as one final fractional step.  There is therefore a fixed $S_b>0$ for
which every stopped state has strict tube margins and
\begin{equation}
 \Gamma\succeq\lambda I/2,\qquad \|\Gamma\|_{\rm op}\le G,\qquad
 \|F\|_{\rm sum}\le V\rho .
 \label{eq:at-pop-tube}
\end{equation}
Here the parameter norm is the sum of the $L^2$ first-layer and readout
norms and the Hilbert--Schmidt hidden-increment norms.  The residual is
also bounded by a fixed constant on this tube.

\paragraph{A response bound uniform in the number of steps.}
For each hidden initialized matrix, introduce its forward sources
$\xi_{a,k}^{(\ell)}$ and reverse sources $\zeta_{a,k}^{(\ell)}$.
Their within-group covariances are respectively
$\E[H_{a,k}^{(\ell-1)}H_{b,r}^{(\ell-1)}]$ and
$\E[\delta_{a,k}^{(\ell)}\delta_{b,r}^{(\ell)}]$.
Expanding the learned operators in \eqref{eq:at-pop-euler} and using
\eqref{eq:at-source-rule} gives
\begin{align}
 Z_{a,k}^{(\ell)}
 &=\xi_{a,k}^{(\ell)}
       +\sum_{b,r<k}\mathsf F_{ak,br}^{(\ell)}\delta_{b,r}^{(\ell)},
 &P_{a,k}^{(\ell-1)}
 &=\zeta_{a,k}^{(\ell)}
       +\sum_{b,r\le k}\mathsf D_{ak,br}^{(\ell)}H_{b,r}^{(\ell-1)},
 \label{eq:at-response-fields}\\
 \mathsf F_{ak,br}^{(\ell)}
 &=\mathsf C_{ak,br}^{(\ell)}
       -\frac{2\alpha_ru_{b,r}}m
          \E[H_{b,r}^{(\ell-1)}H_{a,k}^{(\ell-1)}],
 &\mathsf C_{ak,br}^{(\ell)}
 &=\E\frac{\partial H_{a,k}^{(\ell-1)}}{\partial\zeta_{b,r}^{(\ell)}},
 \nonumber\\
 \mathsf D_{ak,br}^{(\ell)}
 &=\mathsf A_{ak,br}^{(\ell)}
       -\mathbf1_{r<k}\frac{2\alpha_ru_{b,r}}m
          \E[\delta_{b,r}^{(\ell)}\delta_{a,k}^{(\ell)}],
 &\mathsf A_{ak,br}^{(\ell)}
 &=\E\frac{\partial\delta_{a,k}^{(\ell)}}{\partial\xi_{b,r}^{(\ell)}}.
 \nonumber
\end{align}
The root equations are the first-layer and readout recursions in
\eqref{eq:at-pop-euler}.  The training-memory coefficients are bounded
by $J\alpha_r/m$ for a fixed $J$, by the preceding RMS bounds.

For a scalar random variable let
$\mathcal N(U)=\sup_{p\ge2}\|U\|_p/\sqrt p$.
If $\mathcal N(U)\le B$, expansion of the exponential and
$r!\ge(r/e)^r$ give
$\E e^{U^2/(8eB^2)}\le4/3$, and hence
$\E e^{t|U|}\le(4/3)e^{2et^2B^2}$.
Jensen's inequality with weights $\alpha_r/(m\sum\alpha_r)$ therefore
gives, without temporal independence,
\begin{equation}
 \E\exp\left(t\sum_{r<k}\frac{\alpha_r}{m}
                         \sum_b|P_{b,r}^{(\ell)}|\right)
       \le\frac43 e^{2et^2S^2B^2}
 \quad\text{if all }\mathcal N(P_{b,r}^{(\ell)})\le B.
 \label{eq:at-weighted-exponential}
\end{equation}
We verify simultaneously the caps
\begin{equation}
 \sum_{b,r\le k}|\mathsf A_{ak,br}^{(\ell)}|\le\mathfrak a_\ell,
 \qquad
 |\mathsf C_{ak,br}^{(\ell)}|\le\mathfrak c_\ell\alpha_r/m,
 \qquad \mathcal N(H)\le B_H,\quad
 \mathcal N(P^{(\ell)})\le B_{P,\ell}.
 \label{eq:at-response-caps}
\end{equation}
Set $f_\ell=\mathfrak c_\ell+J$ for $\ell\ge2$, $f_1=1$,
$d_\ell=1+\mathfrak a_{\ell+1}$ for $\ell<L$, and $d_L=1$.
Choose $K\ge2$ depending only on the roots, source RMS bounds and
activation bounds, before choosing the response caps.  The field
recursions imply
\begin{align*}
 \mathcal N(H_{a,k}^{(1)})&\le K+KS\sup_{b,r<k}\mathcal N(P_{b,r}^{(1)}),\\
 \mathcal N(H_{a,k}^{(\ell)})&\le K+Kf_\ell S
                              \sup_{b,r<k}\mathcal N(P_{b,r}^{(\ell)}),\\
 \mathcal N(P_{a,k}^{(\ell)})&\le K+(\mathfrak a_{\ell+1}+JS)
                              \sup_{b,r\le k}\mathcal N(H_{b,r}^{(\ell)}),\\
 \mathcal N(w_k)&\le K+KS\sup_{b,r<k}\mathcal N(H_{b,r}^{(L)}).
\end{align*}
The first two inequalities use linear growth, and $\delta=g_\ell(Z)P$
uses only bounded slope.  Thus the forward fields must be bootstrapped
together with the full backward carriers.

Here are the derivative estimates that close this bootstrap.  Let $V_k$
be the maximum through time $k$ of the sum of absolute derivatives of
$Z_a^{(\ell)}$ with respect to its own forward slots.  Differentiating
\eqref{eq:at-response-fields}, the derivative row of $P_a^{(\ell)}$
is bounded by $C d_\ell V_k$; at the top use the readout update instead.
The product rule gives a row bound
$C(|P_a^{(\ell)}|+d_\ell)V_k$ for $\delta_a^{(\ell)}$.  Since the
forward memory uses earlier times only,
\[
 V_k\le1+Cf_\ell\sum_{r<k}\alpha_r
       \left(d_\ell+\frac1m\sum_b|P_{b,r}^{(\ell)}|\right)V_r.
\]
For a single backward source at $(b,r)$ let $D_k$ be the maximal
absolute derivative of the lower preactivation through time $k$.
It is zero for $k\le r$.  At the source time, differentiation of the
direct Gaussian in $P$ gives one pulse; inserting it in the first-layer
update or the next forward memory retains exactly $\alpha_r/m$.
All later derivatives consequently obey
\[
 D_k\le Cf_{\ell-1}\alpha_r/m+
 Cf_{\ell-1}\sum_{u<k}\alpha_u
       \left(d_{\ell-1}+\frac1m\sum_a|P_{a,u}^{(\ell-1)}|\right)D_u.
\]
Discrete Gronwall, \eqref{eq:at-weighted-exponential}, and
Cauchy--Schwarz for the terminal carrier factor yield
\begin{align}
 \sum_{b,r\le k}|\mathsf A_{ak,br}^{(\ell)}|
 &\le C(B_{P,\ell}+d_\ell)
       e^{Cf_\ell Sd_\ell+Cf_\ell^2S^2B_{P,\ell}^2},\nonumber\\
 \frac{|\mathsf C_{ak,br}^{(\ell)}|}{\alpha_r/m}
 &\le Cf_{\ell-1}
       e^{Cf_{\ell-1}Sd_{\ell-1}
                          +Cf_{\ell-1}^2S^2B_{P,\ell-1}^2}.
 \label{eq:at-cap-improvement}
\end{align}
Fix one sufficiently large $C$ in these inequalities independently of
the caps.  Choose, in order,
\[
 \mathfrak c_2=4C,\qquad
 \mathfrak c_\ell=4C(\mathfrak c_{\ell-1}+J)\quad(\ell>2),
 \qquad
 \mathfrak a_L=4C(2K+1),\qquad
 \mathfrak a_\ell=4C(4K+1)(1+\mathfrak a_{\ell+1}) .
\]
The $\mathfrak c$ caps are selected from bottom to top, the
$\mathfrak a$ caps from top to bottom.  Set
$B_H=B_{P,L}=2K$ and
$B_{P,\ell}=4K(1+\mathfrak a_{\ell+1})$ for $\ell<L$.
All these constants are now fixed.  Choose $S_r>0$ so $JS_r\le1$,
$2KS_r\le1$, $f_\ell S_rB_{P,\ell}\le1$, and every exponent in
\eqref{eq:at-cap-improvement} is at most $\log2$.
The field recursions preserve their caps, and the response estimates
improve theirs by a factor of two.

This is a chronological induction, not an assumption on future fields.
At node $k$, the first layer and readout use only earlier nodes.
Construct current forwards from bottom to top: each current
$\mathsf C^{(\ell)}$ uses the already constructed lower field and past
lower carriers, after which $H^{(\ell)}$ uses only past same-layer
carriers.  Construct current backwards from top to bottom: the current
upper response supplies the next carrier, and its derivative row supplies
the next $\mathsf A$.  At $k=0$ forward memory is empty and the readout
is the zero root.  This starts the induction.  A fractional last Euler
step changes only its nonnegative activity weight.  We have proved, for
every program stopped at $S\le\min(S_b,S_r)$, including every affine
interpolation state,
\begin{equation}
 \max_{\ell,a}\E e^{c|U_a^{(\ell)}|^2}\le C,
 \qquad U\in\{Z,H,P,\delta\}.
 \label{eq:at-pop-gaussian}
\end{equation}
The constants are independent of physical duration and number of steps.

\paragraph{Excluding the stop and constructing the strong flow.}
On the parameter tube, forward subtraction costs $Cd$ in the parameter
distance.  For a reference carrier $P_*$, the only difficult backward
product is controlled by
\begin{equation}
 \|[g_\ell(Z)-g_\ell(Z_*)]P_*\|_2
 \le jR\|Z-Z_*\|_2+2s\|P_*\mathbf1_{\{|P_*|>R\}}\|_2 .
 \label{eq:at-reference-cutoff}
\end{equation}
The remaining backward difference is propagated by bounded operators
and bounded gates, so new cutoff contributions add through fixed depth.
The rank-one identity and \eqref{eq:at-pop-gaussian} therefore give,
for the prediction differential $J$ from the mobility Hilbert space to
sample RMS,
\begin{equation}
 \|J(\vartheta)-J(\theta_*)\|
 \le C\{(1+R)d(\vartheta,\theta_*)+e^{-cR^2}\}.
 \label{eq:at-jacobian-modulus}
\end{equation}
Only the reference node needs tails.  The scalar prediction is
continuously differentiable in these parameter spaces.  To check the
possibly unbounded multiplier explicitly, its scalar Taylor remainder
satisfies, for a fixed $P\in L^2$,
\[
 |\E[P\{\phi(Z+h)-\phi(Z)-g(Z)h\}]|
 \le \tfrac12jR\|h\|_2^2+
       2s\|P\mathbf1_{\{|P|>R\}}\|_2\|h\|_2=o(\|h\|_2).
\]
First send $h$ to zero at fixed $R$, then $R$ to infinity.
Forward subtraction and descending this weighted scalar expansion give
the displayed Jacobian; bounded-multiplier continuity proves its
continuity.  This does not assume that the full activation map is
Fréchet differentiable from $L^2$ to $L^2$.

An Euler segment has parameter length at most $Vh_k\rho_k$.
Integrating the prediction differential on that segment and using
\eqref{eq:at-jacobian-modulus} gives
\[
 r_{k+1}=(I-2h_k\Gamma_k)r_k+e_k,\qquad
 \|e_k\|_m\le h_k\rho_k\omega(h_k),\qquad
 \omega(h)=C\inf_{R\ge1}\{(1+R)h+e^{-cR^2}\}\longrightarrow0.
\]
For sufficiently small maximal physical step,
\eqref{eq:at-pop-tube} implies
\[
 \rho_{k+1}\le(1-\lambda h_k/2)\rho_k,\qquad
 \sum_{k<j}h_k\rho_k\le2Y/\lambda.
\]
Choose $S=4Y/\lambda$ and reduce $Y_*$ so
$S\le\min(S_b,S_r)$.
The first attainment of this activity cap, even by a shortened final
step, contradicts the strict factor-two margin.  The physical tube
already has a strict margin.  These deterministic Euler programs
therefore extend through every prescribed physical horizon.

Put countably many refining programs on all integer horizons on the
common generated spaces above.  The vector-field version of
\eqref{eq:at-jacobian-modulus} follows from backward subtraction and
rank-one update subtraction.  Comparing two Euler interpolants on
$[0,T]$, with maximal steps $h,h'$, gives
\begin{equation}
 \sup_{t\le T}d(\theta^h(t),\theta^{h'}(t))
 \le C_Te^{C(1+R)T}
       \{(1+R)(h+h')+e^{-cR^2}\}.
 \label{eq:at-pop-cauchy}
\end{equation}
The same argument holds in the Hilbert--Schmidt increment norm because
$\|u\otimes v-u'\otimes v'\|_{\rm HS}
\le\|u-u'\|_2\|v\|_2+\|u'\|_2\|v-v'\|_2$.
First send the meshes to zero at fixed $R$, then $R$ to infinity.
Since $e^{CRT-cR^2}\to0$, completeness gives a uniform strong limit.
For the passage through backward gates, if $Z_i\to Z$ in probability
and $P_i\to P$ in $L^2$, split
$g(Z_i)P_i-g(Z)P$ into $g(Z_i)(P_i-P)$ and
$[g(Z_i)-g(Z)]P$.  The former converges by boundedness of $g$;
truncate the single integrable variable $P^2$ to prove convergence of
the latter.  Thus the vector field is continuous, and the Euler
integral equations pass to the limit.

The same reference comparison against this constructed solution proves
uniqueness: the discrepancy of two solutions is at most
$C_Te^{CRT-cR^2}$ for every $R$, hence zero, up to any first tube exit.
Equality to the constructed path excludes that exit.  Solutions from
different horizons agree on overlaps.  The Gram gap survives the limit,
and the exact residual equation $\dot r=-2\Gamma r$ gives exponential
decay and finite total parameter variation, with the constants in
\eqref{eq:at-fitting}--\eqref{eq:at-speeds} decreased or enlarged if
necessary.  The limit at infinite physical time is an interpolating
parameter state.  At each time, $L^2$ convergence and Fatou pass
\eqref{eq:at-pop-gaussian} to the solution; the constants are common to
all times.  Applying the same reasoning at the parameter limit includes
$t=\infty$.  The assertion bounds time marginals; it places no time
supremum inside the exponential.

Finally remove classical second differentiability.  Compactly supported
mollification gives
$\|\phi_{\ell,\epsilon}-\phi_\ell\|_\infty\le Cs\epsilon$,
$\|g_{\ell,\epsilon}-g_\ell\|_\infty\le Cj\epsilon$,
$\|g_{\ell,\epsilon}\|_\infty\le s$ and
$\|g_{\ell,\epsilon}'\|_\infty\le j$.
The value at zero is uniformly bounded.  All response caps and activity
thresholds are therefore common to sufficiently small smoothings.
The initialization recursion \eqref{eq:at-initial-covariance} is continuous
under these uniform activation perturbations, so its positive Gram margin
also persists, after decreasing the common margin once.
The comparison \eqref{eq:at-pop-cauchy} between different smoothings
has the additional term $C(1+R)(\epsilon+\epsilon')$ inside braces.
After constructing each smoothed flow, take smoothing scales to zero at
fixed $R$ and then $R$ to infinity.  This gives the strong solution for
the stated $C^1$ activations with Lipschitz derivatives, with the same
uniqueness and Gaussian bounds.

\begin{lemma}[Dense carrier transfer]
\label{lem:at-gaussian-tail}
On the events \eqref{eq:at-good-event}, let
$k_{a,D}^{(L)}=w_D$ and
$k_{a,D}^{(\ell)}=(W_D^{(\ell+1)})^\top\delta_{a,D}^{(\ell+1)}$
for $\ell<L$.  Define
\[
 H_n(M,t)=\sum_{\ell=1}^L\max_a
       \frac{\|k_{a,D}^{(\ell)}(t)
           \mathbf1_{\{|k_{a,D}^{(\ell)}(t)|>M\}}\|_2}{\sqrt n},
 \qquad Z_n(M)=\int_0^\infty\rho_D(t)H_n(M,t)\dd t,
\]
and set $Z_n(M)=0$ off $\mathcal G_n$.  There are constants $c,C>0$
and dense-only random variables $a_n\ge0$, with $a_n\to0$ in
probability, such that, simultaneously for all integer $M\ge1$,
\begin{equation}
 Z_n(M)\le Ce^{-cM^2}+a_n.
 \label{eq:at-dense-tail}
\end{equation}
On each finite physical horizon, dense finite-width predictions on any
fixed finite list of training or passive test inputs converge in probability,
uniformly in time, to the predictions of the unique strong population
flow just constructed.
\end{lemma}
\begin{proof}
Fix a finite horizon $T$, a reference cutoff $R$, and a finite population
Euler mesh.  Evaluate its expanded initialized-matrix calls on the actual
finite Gaussian arrays, keeping every residual and scalar contraction
equal to its deterministic population value.  This is a fixed finite
program to which Lemma~\ref{lem:at-fixed-program} applies.
Reconstruct finite proxy parameters by the rank-one updates with those
oracle vectors.  For example a proxy action on an oracle feature differs
from its prescribed action by a finite sum of terms of the form
\[
 -\frac{2\alpha_ru_{b,r}}m\,\delta_{b,r}^{(\ell)}
 \left\{\frac{(h_{b,r}^{(\ell-1)})^\top h_{a,k}^{(\ell-1)}}n
              -\E[H_{b,r}^{(\ell-1)}H_{a,k}^{(\ell-1)}]\right\}.
\]
Each brace converges to zero; all involved vector RMS norms are bounded
in probability.  Forward recomputation is therefore consistent.
For backward recomputation use \eqref{eq:at-reference-cutoff} against
the oracle carriers.  Continuous quadratic-growth cutoff measurements
of their joint empirical laws converge by the lemma.  At fixed cutoff
the consistency errors vanish, and the limiting cutoff tails are
$Ce^{-cR^2}$ by \eqref{eq:at-pop-gaussian}.  Sending that cutoff to
infinity proves RMS consistency of the recomputed backward fields,
residuals, and vector field.

Compare actual finite dense gradient flow to this proxy interpolant,
stopping it only if it leaves a slightly larger physical tube.
The identical reference inequality and integral comparison give
\[
 \sup_{t\le T}d_n(\theta_{n,D}(t),\theta_{n,\mathrm{proxy}}^h(t))
 \le C_Te^{C(1+R)T}
       \{(1+R)h+e^{-cR^2}+o_{\mathbb P}(1)\}.
\]
The last term is at fixed $h,R$.  The proof uses empirical tails only
of fixed-program reference nodes.  Taking width to infinity at fixed
$h,R$, then $h$ to zero, then $R$ to infinity makes the error smaller
than the exit margin and excludes the stop.  Combining this with
\eqref{eq:at-pop-cauchy} proves the compact-time observable convergence.
For the original $C^{1,1}$ activations insert a fixed smooth reference
program and its $C(1+R)\epsilon$ comparison error, taking its smoothing
scale to zero before the final $R$ limit.  Passive test inputs require
only finitely many extra forward queries and first-layer root coordinates;
they change no training update.  They are therefore covered by exactly
the same fixed-program and comparison arguments.

We spell out how this yields actual dense carrier tails.  For fixed
$M$, use the continuous 1-Lipschitz function
$g_M(v)=(|v|-M)_+$, with
$|v|\mathbf1_{\{|v|>2M\}}\le2g_M(v)$.
Its normalized Euclidean norm changes by at most the RMS change in
the carrier.  Fixed-program convergence transfers its squared norm at
every reference node.  A finite time grid, the uniform parameter-speed
bound, and \eqref{eq:at-reference-cutoff} against these reference nodes
control the interpolation error.  First choose the grid and a comparison
cutoff, let width tend to infinity at that fixed program, then refine
the grid and remove the comparison cutoff.  The population Gaussian
bound implies a limiting carrier cutoff norm at most $Ce^{-cM^2}$,
after halving $M$ and changing $c$.  Prediction convergence also
transfers the residual factor.  We conclude that for every fixed
$M\ge1$ and $\epsilon>0$,
\[
 \Pr\left\{\int_0^T\rho_D(t)H_n(M,t)\dd t
                   >Ce^{-cM^2}+\epsilon\right\}\longrightarrow0.
\]
The constant is independent of $T$: population carrier tails are uniform
in time and total residual activity is bounded.
On $\mathcal G_n$, finite dense RMS bounds imply $H_n(M,t)\le C$,
and \eqref{eq:at-fitting} gives
$\int_T^\infty\rho_DH_n\dd t\le Ce^{-\kappa T}$
uniformly in $n,M$.  Choosing $T$ after the error tolerance but before
the width limit proves
\[
 \Pr\{Z_n(M)>Ce^{-cM^2}+\epsilon\}\longrightarrow0
 \quad\text{for every fixed }M\ge1,\ \epsilon>0.
\]
Set
\[
 a_n=\sup_{M\in\mathbb N,\ M\ge1}
                  \bigl(Z_n(M)-Ce^{-cM^2}\bigr)_+.
\]
These random variables are measurable, bounded, and depend only on the
dense flow and its initialization.  To prove $a_n\to0$, fix a tolerance
and choose one integer $J$ with $Ce^{-cJ^2}$ smaller than half that
tolerance.  Monotonicity of $Z_n$ bounds every excess at $M\ge J$
by $Z_n(J)$, which is small in probability.  The finitely many smaller
integers are handled by a union bound and fixed-cutoff convergence.
This proves \eqref{eq:at-dense-tail} simultaneously in $M$, with no
width-dependent cutoff theorem and no numerical width rate.
\end{proof}
